% Date: May 10, 2026
% Status: Resubmission - Major Revision (Stabilized L=8, K=2)
%==============================================================================

\documentclass[journal=jpcl,manuscript=letter]{achemso}
\usepackage[utf8]{inputenc}
\usepackage[T1]{fontenc}
\usepackage{amsmath,amssymb,graphicx,booktabs,siunitx,physics}
\usepackage{xcolor}
\usepackage{xr-hyper}
\usepackage[colorlinks=true,linkcolor=blue,citecolor=blue,urlcolor=blue]{hyperref}
\usepackage{cleveref}


\externaldocument[SI-]{SI_JPCL_26-05-10}

% SIunitx configuration
\sisetup{
    locale = US,
    detect-all = true,
    group-minimum-digits = 4,
    separate-uncertainty = true,
    multi-part-units = single,
    range-phrase = {--},
    range-units = single,
    number-unit-product = \,,
    input-comparators = { < > = \le \ge \ll \gg \leq \geq }
}

%==============================================================================
% AUTHOR INFORMATION
%==============================================================================

\title{Quantum-coherent spectral engineering: non-Markovian dynamics in photosynthetic energy transfer under engineered photon baths}

\author{Steve Cabrel Teguia Kouam}
\affiliation{Department of Physics, Faculty of Science, University of Douala, Cameroon}
\email{steve.teguia@facsciences-uy1.cm}

\author{Theodore Goumai Vedekoi}
\affiliation{Department of Physics, Faculty of Science, University of Yaoundé I, Cameroon}

\author{Jean-Pierre Tchapet Njafa}
\affiliation{Department of Physics, Faculty of Science, University of Yaoundé I, Cameroon}

\author{Jean-Pierre Nguenang}
\affiliation{Department of Physics, Faculty of Science, University of Douala, Cameroon}

\author{Serge Guy Nana Engo}
\affiliation{Department of Physics, Faculty of Science, University of Yaoundé I, Cameroon}

%==============================================================================
% DOCUMENT BEGIN
%==============================================================================

\begin{tocentry}
\includegraphics[width=\textwidth]{TOC_Graphic.png}
\end{tocentry}

\begin{document}

\begin{abstract}
By partitioning the photonic bath into resonant and off-resonant modes, we demonstrate a mechanism for \textbf{dephasing suppression} in photosynthetic energy transfer. Aligning the pump spectrum with underdamped vibronic resonances in the Fenna--Matthews--Olson (FMO) complex prepares polaron-dressed states effectively decoupled from dissipative fluctuations, extending excitonic coherence lifetimes by \SIrange{20}{50}{\percent} and improving forward transfer yields by \SI{25}{\percent} at \SI{295}{\kelvin}. Using process tensor hierarchy of pure states (PT-HOPS/SBD), we show that dual-band filtering at \SIlist{750;820}{\nano\meter} selectively targets vibronic resonances while bypassing dephasing-dominant noise. This enhancement is robust against static disorder ($\sigma = \SI{50}{\per\centi\meter}$), with an ensemble-averaged increase of $\eta = 0.20 \pm 0.04$. These results establish \textbf{spectral bath engineering} as a foundational design principle for coherence-assisted transport. This framework extends to \textbf{symbiotic agrivoltaic systems}, where organic photovoltaics act as active spectral filters to co-optimize excitonic transport alongside the photosynthetic requirements of underlying crops.
\end{abstract}

The observation of long-lived quantum coherence in biological light-harvesting complexes has remained a central theme in molecular biophysics since two-dimensional electronic spectroscopy (2DES) first revealed oscillatory signals in the Fenna--Matthews--Olson (FMO) complex.\cite{Engel2007,Panitchayangkoon2010} This sparked a vigorous debate regarding the functional relevance of quantum mechanical effects in warm, noisy biological environments.\cite{Cao2020,Ishizaki2009a} While initial interpretations attributed these signals to purely electronic coherences, subsequent investigations highlighted the critical role of vibronic coupling, where electronic transitions are strongly mixed with underdamped vibrational modes of the protein scaffold.\cite{Christensson2012,Chin2013,Fleming2015,Scholes2015}

At physiological temperatures (\SI{295}{\kelvin}), the decoherence of electronic superpositions is typically expected to occur on sub-\SI{100}{\femto\second} timescales due to the dense manifold of thermal fluctuations in the solvent-protein environment. However, experimental evidence suggests that coherence lifetimes in FMO can extend to $\gtrsim \SI{300}{\femto\second}$, a duration comparable to the primary energy transfer steps.\cite{Engel2007,Hildner2013} This surprising longevity is often attributed to the non-Markovian nature of the environment, where the bath retains a memory of past system states, thereby slowing the loss of quantum information.\cite{Tanimura1989,Suess2014}

Existing theoretical frameworks describe how these environments might optimize energy transport. These include environment-assisted quantum transport (ENAQT), where noise assists navigation of the excitonic landscape,\cite{Mohseni2008,Rebentrost2009NJP,Plenio2010NJP,Wu2010NewJPhys} and vibronically-mediated transport, where vibrational resonances bridge energy gaps between chromophores.\cite{Chin2013,Huelga2013} While the spectral environment is usually treated as an immutable property, whether the effective system-bath interaction can be co-optimized with environmental factors to navigate the parameter space toward maximal energy transfer efficiency remains an open question.

We propose \textbf{quantum control via selective vibronic excitation}: shaping an incident impulsive laser pulse to selectively populate vibronic resonances that protect excitonic coherence. Rather than attempting to suppress all environmental interactions, this approach partitions the bath into protected resonant modes and decoherent off-resonant modes. Using process tensor hierarchy of pure states (PT-HOPS),\cite{Citty2024a} we find that aligning the pump spectrum with vibronic peaks in FMO (\SIlist{750;820}{\nano\meter}) extends coherence lifetimes.

Dual-band spectral filtering leads to a \SIrange{20}{50}{\percent} increase in coherence lifetime and a \SI{25}{\percent} enhancement in forward transfer yield at room temperature. This effect persists under physiological disorder and temperature fluctuations. This work addresses the \textbf{energy--water--food nexus} by providing a quantum-mechanical basis for agrivoltaics, where the photonic bath is no longer an immutable source of noise but a controllable resource for symbiotic energy-food production. This "quantum-to-organism" perspective\cite{Vandamme2025} allows us to link microscopic coherence to macroscopic energy yields. Below, we outline a protocol for verifying these results in 2DES experiments using spatial light modulator (SLM) pulse-shaping, suggesting selective vibronic excitation as a strategy for enhancing quantum transport in molecular materials.

\begin{figure}[htpb]
\centering
\includegraphics[width=0.8\columnwidth]{FMO_Schematic_JPCL.png}
\caption{\textbf{Schematic of the Fenna--Matthews--Olson (FMO) complex and spectral engineering.} The seven-site bacteriochlorophyll $a$ network is driven by a dual-band impulsive pulse. The incident field is shaped by a spectral filter to selectively excite underdamped vibronic resonances, promoting enhanced coherence and directed energy transfer toward the reaction center.}
\label{fig:fmo_schematic}
\end{figure}

To investigate the dynamics of excitonic energy transfer under engineered photon baths, we employ the Process Tensor Hierarchy of Pure States (PT-HOPS) method,\cite{Suess2014,Citty2024a} which provides a numerically exact treatment of non-Markovian open quantum systems. PT-HOPS is the theoretical non-perturbative formalism used to handle non-Markovian memory without exponential scaling; all simulations are performed using the open-source \textsc{MesoHOPS} library (v1.6),\cite{Varvelo2021a} which implements PT-HOPS with Stochastically Bundled Dissipators (SBD) for size-invariant $\mathcal{O}(1)$ scaling with system size. The FMO complex is modeled as a seven-site excitonic system where each chromophore (bacteriochlorophyll $a$) is coupled to its own independent phonon bath.

The total Hamiltonian of the system and its environment is given by:
\begin{equation}
\hat{H} = \hat{H}_S + \hat{H}_B + \hat{H}_{SB} + \hat{H}_{ext}(t),
\label{eq:total_hamiltonian}
\end{equation}
where $\hat{H}_S$ is the excitonic Hamiltonian of the protein-pigment complex:
\begin{equation}
\hat{H}_S = \sum_{n=1}^{7} \varepsilon_n \dyad{n} + \sum_{n \neq m} J_{nm} \dyad{n}{m}.
\label{eq:excitonic_hamiltonian}
\end{equation}
Here, $\varepsilon_n$ denotes the site energy of the $n$-th pigment, and $J_{nm}$ represents the electronic coupling between sites $n$ and $m$. The environmental interaction is modeled by $\hat{H}_{SB} = \sum_n \hat{L}_n \otimes \hat{B}_n$, where $\hat{L}_n = \dyad{n}$ are the coupling operators and $\hat{B}_n$ represent the collective coordinates of the $n$-th bath.

The bath's influence is fully characterized by the 12-mode Kleinekathöfer/Coker spectral density $J(\omega)$, which we decompose into a continuous, overdamped Drude--Lorentz component at \SI{295}{\kelvin} and a set of underdamped vibronic modes (\cref{fig:quantum_dynamics}e):
\begin{equation}
J(\omega) = 2\lambda_D \frac{\omega \gamma_D}{\omega^2 + \gamma_D^2} + \sum_{k} 2\lambda_k \frac{\omega \omega_k^2 \gamma_k}{(\omega_k^2 - \omega^2)^2 + \omega^2 \gamma_k^2}.
\label{eq:spectral_density}
\end{equation}
The first term accounts for the broad protein-solvent fluctuations with reorganization energy $\lambda_D$ and cutoff $\gamma_D$. The second term represents specific underdamped vibrational modes with frequencies $\omega_k$, reorganization energies $\lambda_k$, and damping rates $\gamma_k$.

The PT-HOPS method evolves a hierarchy of pure states $\ket{\psi^{(\vb{n})}(t)}$, indexed by the multi-index $\vb{n}$, which captures the system-bath correlations to arbitrary order. To model the spectrally shaped impulsive pump, the initial state is constructed via an exact physical projection of the transmission-filtered field, $E_{\mathrm{eff}}(\omega) = T(\omega) E_{\mathrm{in}}(\omega)$, directly onto the excitonic eigenvectors. The temporal envelope is defined as a Gaussian $E(t) = E_0 \exp(-t^2/2\sigma_t^2)$ with a full-width at half-maximum (FWHM) of \SI{50}{\femto\second} centered at $t = 0$. Here, $T(\omega)$ is the transmission function of the engineering filter:
\begin{equation}
T(\omega) = \sum_{j=1}^2 \exp\left[ -\frac{(\omega - \Omega_{j})^2}{2\Delta\omega^2} \right].
\label{eq:transmission_function}
\end{equation}
We target the primary vibronic resonances at $\Omega_1 = \SI{750}{\nano\meter}$ and $\Omega_2 = \SI{820}{\nano\meter}$ with a bandwidth $\Delta\omega = \SI{100}{\per\centi\meter}$. %==============================================================================
% FIGURES: GRAPHICAL ABSTRACT
%==============================================================================

\begin{figure}[h]
\centering
\includegraphics[width=\textwidth]{graphical_abstract.png}
\caption{\textbf{Graphical Abstract}: Workflow for quantum-enhanced agrivoltaics via spectral bath engineering. Solar irradiance is spectrally filtered targeting vibronic resonances in the FMO complex. PT-HOPS simulations quantify coherence extension and transfer yield enhancement, informing sustainable energy harvest in agricultural environments.}
\label{fig:graphical_abstract}
\end{figure}

The hierarchy is truncated at $L=8$ with $K=2$ Matsubara terms and a time step $\Delta t = \SI{1.0}{\femto\second}$, providing rigorously converged results at \SI{295}{\kelvin} (truncation MAE $\approx 3.42 \times 10^{-5}$) while ensuring memory stability on mesoscale hardware (128 GB). We employ an integrated resource-aware architecture (\cref{fig:resource_architecture}) that dynamically gates parallel worker threads based on available RAM and heuristic memory estimation per trajectory, ensuring numerical stability and preventing Out-Of-Memory (OOM) failures across both high-performance clusters and standard workstation environments. Numerical validations against the Hierarchical Equations of Motion (HEOM) and physical consistency checks are detailed in SI Sections S1--S4. The PT-HOPS/SBD framework was verified for hardware-agnostic reproducibility, maintaining 100% trace preservation and identical excitonic populations across heterogeneous compute environments.

\begin{figure}[h]
\centering
\includegraphics[width=0.8\columnwidth]{resource_architecture.png}
\caption{\textbf{Resource-Aware Simulation Architecture}: Flowchart showing trajectory parameter input, dynamic heuristic memory estimation, real-time system RAM monitoring, and autonomous thread gating (66.7\% limit) for stable parallel execution.}
\label{fig:resource_architecture}
\end{figure}

%==============================================================================
% RESULTS AND DISCUSSION
%==============================================================================

\textbf{Results and Discussion.}

\subsection*{Coherence Extension}

\begin{figure}[t]
\centering
\begin{minipage}{0.65\columnwidth}
  \includegraphics[width=\textwidth]{Quantum_dynamics.png}
\end{minipage}
\hfill
\begin{minipage}{0.32\columnwidth}
  \includegraphics[width=\textwidth]{Figure1e_spectral_relationships.pdf}
\end{minipage}
\caption{\textbf{Quantum dynamics and spectral relationships.} PT-HOPS/SBD simulations at \SI{295}{\kelvin} with static disorder $\sigma = \SI{50}{\per\centi\meter}$ (\num{100} disorder realizations). \textbf{(a)}~$l_1$-norm coherence $C_{l_1}(t)$ showing \SIrange{20}{50}{\percent} lifetime extension under dual-band filtering. Solid and dashed lines denote broadband and filtered excitation, respectively. Insets display the relative difference. \textbf{(b)}~Site populations demonstrating \SI{25}{\percent} enhancement in forward transfer yield. \textbf{(c)}~Inverse participation ratio (IPR) quantifying exciton delocalization. \textbf{(d)}~Quantum Fisher Information (QFI). Shaded regions: \SI{95}{\percent} confidence intervals. \textbf{(e)}~Spectral relationships: FMO absorption (green), incident broadband pulse $E_{\mathrm{in}}$ (orange), bath spectral density $J_{\mathrm{bath}}$ (blue), and dual-band filter transmission $T(\omega)$ (red peaks at \SIlist{750;820}{\nano\meter}).}
\label{fig:quantum_dynamics}
\end{figure}

We quantify coherence using the $l_1$-norm,\cite{Baumgratz2014} a basis-dependent measure accessible through state tomography:
\begin{equation}
C_{l_1}(\rho) = \sum_{i \neq j} \abs{\rho_{ij}}.
\label{eq:l1_coherence}
\end{equation}
We define the coherence lifetime $\tau_c$ as the $1/e$ decay time of $C_{l_1}(t)$. Under optimal dual-band filtering, $\tau_c = \SI{420(35)}{\femto\second}$ at \SI{295}{\kelvin}, compared with $\tau_c = \SI{280(25)}{\femto\second}$ under broadband impulsive excitation---corresponding to a \SI{50}{\percent} extension at the upper bound of the confidence interval (\cref{fig:quantum_dynamics}a).

The Inverse Participation Ratio (IPR), $\mathrm{IPR} = (\sum_n \abs{\mel{n}{\hat{\rho}}{n}}^2)^{-1}$, quantifies exciton delocalization across the chromophore network. Under broadband impulsive excitation, the IPR decays rapidly from its initial peak, settling at a time-averaged value of $\mathrm{IPR} \approx 4.1$. In contrast, spectral filtering maintains a higher degree of delocalization, with $\mathrm{IPR} \approx 6.8$ (filtered), as shown in \cref{fig:quantum_dynamics}c. 

This enhancement in delocalization results from reduced system-bath entanglement in the polaron frame. By selectively populating states effectively decoupled from the most dissipative bath modes, the system preserves its wave-like character. In the FMO complex, the enhanced IPR (quantified by $\mathrm{IPR} \approx 6.8$) enables the exciton to sample a broader spatial manifold of the chromophore network. This facilitates \textbf{constructive interference among multiple transfer pathways}, allowing the system to effectively bypass energetic traps that would otherwise lead to localization and recombination in a classical hopping regime.

\textbf{Forward transfer yield enhancement.}
We evaluate the forward transfer yield ($\Phi_{\mathrm{FT}}$), defined as the fractional depletion of the initially excited antenna state:
\begin{equation}
\Phi_{\mathrm{FT}} = 1 - \Tr[\hat{P}_{1}\,\vb*{\rho}(t_{\max})],
\label{eq:ETE}
\end{equation}
where $\hat{P}_{1} = \dyad{1}$ projects onto the initial excitation site (BChl 1 or 6) and $t_{\max} = \SI{1}{\pico\second}$. Spectral filtering enhances $\Phi_{\mathrm{FT}}$ by \SI{25(3)}{\percent} relative to broadband excitation (\cref{tab:quantum_metrics}). Statistical significance was rigorously controlled via ensemble averaging; the sampling error of the mean enhancement is bounded at $\pm \SI{3.1}{\percent}$ ($N=100$), ensuring the reported forward transfer increase is robust against stochastic noise. The high-rigor $N=100$ ensemble was enabled by an OOM-hardened parallelization pipeline, which utilized a dynamic memory-aware job scheduler to maintain absolute trace preservation on mesoscale hardware. This metric serves as a robust proxy for the energy available to the reaction center under trace-preserving dynamics.

The enhancement mechanism involves selective population of vibronically resonant states. The filtered driving field preferentially excites modes satisfying:
\begin{equation}
\Delta E_{nm} \approx \hbar\omega_{\mathrm{vib}} \pm J_{nm},
\label{eq:resonance_condition}
\end{equation}
where $\Delta E_{nm} = \varepsilon_n - \varepsilon_m$ is the inter-site energy gap, $\omega_{\mathrm{vib}}$ is the underdamped vibronic frequency, and $J_{nm}$ is the electronic coupling. For the dominant $\ket{1} \to \ket{3}$ pathway ($\Delta E \approx \SI{150}{\per\centi\meter}$), the $\omega_{\mathrm{vib}} = \SI{180}{\per\centi\meter}$ mode (Vibronic Mode 1) mediates resonance.\cite{Chin2013,Fleming2015,Scholes2015}

The 59\% increase in Quantum Fisher Information (QFI) supports the effectiveness of this approach. QFI, $F_Q[\rho, \hat{O}]$, quantifies the information a state $\rho$ contains about a generator $\hat{O}$ (e.g., the excitonic Hamiltonian):
\begin{equation}
F_Q[\rho, \hat{O}] = \sum_{\lambda_i + \lambda_j > 0} \frac{2}{(\lambda_i + \lambda_j)} \abs{\mel{i}{[\hat{\rho}, \hat{O}]}{j}}^2,
\label{eq:QFI}
\end{equation}
where $\lambda_i, \ket{i}$ are the eigenvalues and eigenvectors of $\rho$. In light-harvesting, higher QFI indicates that the excitonic state maintains purity and operational coherence despite the environment.\cite{Pezze2009} The results suggest that spectral filtering prepares states resilient to the information-scrambling effects of the phonon bath.

\textbf{Physical mechanism.}
The coherence and yield enhancements stem from a two-stage process involving both state preparation and subsequent decoherence dynamics. 

\textit{Stage 1: Selective vibronic driving.} The spectral filter $T(\omega)$ aligns the incident field's intensity with the underdamped vibronic modes of the FMO complex (\cref{fig:quantum_dynamics}e). Specifically, the peak at \SI{750}{\nano\meter} (\SI{13333}{\per\centi\meter}) targets the BChl 1 and 6 antenna states, which are strongly coupled to the reaction center through middle-energy exciton states. By exciting these antenna sites with a field that is spectrally tuned to the $\SI{180}{\per\centi\meter}$ and $\SI{575}{\per\centi\meter}$ vibrational frequencies, we prepare polaron-dressed states rather than bare excitonic superpositions. These dressed states are the "transport-optimized" eigenfunctions of the coupled system-bath Hamiltonian.\cite{Chin2013} The primary advantage of preparing these dressed states lies in their decoherence properties. In the standard excitonic basis, pure dephasing is driven by the entire spectral density $J(\omega)$. However, by transforming to a polaron frame (see SI Section~S1.6 for the mathematical derivation), the effective system-bath coupling is redefined. The part of the bath that is "correlated" with the system (the resonant vibronic modes) becomes part of the dressed system, leaving only the residual, off-resonant bath to cause decoherence. Our results show that this "residual" reorganization energy $\lambda_{\mathrm{res}}$ is significantly lower than the total reorganization energy of the broadband case. For the FMO parameters, $\lambda_{\mathrm{res}}$ is reduced by \SI{25}{\percent}, which translates to a \SI{50}{\percent} extension in the $1/e$ coherence lifetime, as the pure dephasing rate in the Markovian limit scales as $\tau_c \propto 1/\lambda_{\mathrm{res}}$.

\textbf{Quantum Metrics Summary.}
\cref{tab:quantum_metrics} summarizes the quantum metric enhancements. The \SI{89}{\percent} concurrence enhancement reflects strengthened inter-site entanglement, while the \SI{59}{\percent} increase in Quantum Fisher Information (QFI) witnesses quantum-enhanced sensitivity.

\begin{table}[t]
\centering
\caption{\textbf{Quantum metric enhancement under spectral filtering.} Dual-band (\SIlist{750;820}{\nano\meter}) vs.\ broadband impulsive excitation at \SI{295}{\kelvin}. Ensemble averages over \num{100} static disorder realizations ($\sigma = \SI{50}{\per\centi\meter}$); uncertainties denote \SI{95}{\percent} confidence intervals. $\Phi_{\mathrm{FT}}$ values expressed as percentages of maximal theoretical yield (\SI{100}{\percent} = unity quantum yield).}
\label{tab:quantum_metrics}
\begin{tabular}{lccc}
\toprule
\textbf{Metric} & \textbf{Filtered} & \textbf{Broadband} & \textbf{Relative Enhancement} \\
\midrule
$\Phi_{\mathrm{FT}}$ (\si{\percent}) & 89.2(21) & 71.4(18) & $0.25 \pm 0.03$ \\
$\tau_c$ (\si{\femto\second}) & 420(35) & 280(25) & $0.50 \pm 0.18$ \\
IPR (sites) & 6.8(5) & 4.1(4) & $0.66 \pm 0.15$ \\
QFI (max) & 12.4(11) & 7.8(8) & $0.59 \pm 0.17$ \\
Concurrence & 0.34(5) & 0.18(4) & $0.89 \pm 0.42$ \\
\bottomrule
\end{tabular}
\end{table}

\textbf{Generalizability.}
To demonstrate that quantum control via selective vibronic excitation is a general property of open quantum systems satisfying the resonance condition (Eq.~\cref{eq:resonance_condition}), we applied the identical filtering protocol to a minimal three-site excitonic trimer. This trimer model, detailed in SI Section~S10, features energy gaps and vibronic couplings tuned to mimic different biological and synthetic systems. The sites represent a generic donor-bridge-acceptor configuration, with the donor and acceptor coupled to an underdamped vibronic mode of frequency $\omega_{\mathrm{vib}} = \SI{200}{\per\centi\meter}$ and reorganization energy $\lambda_{\mathrm{vib}} = \SI{10}{\per\centi\meter}$.

Dual-band filtering robustly yields a fractional enhancement of $\eta_{\mathrm{FT}} = 0.18 \pm 0.03$ and a coherence lifetime extension of \SI{35}{\percent} for the trimer. Critically, we find that the enhancement is highly sensitive to the spectral overlap between the filter peaks and the vibronic sidebands of the exciton states. An off-resonant filter (detuned by \SI{200}{\per\centi\meter}) produces an enhancement of $\eta < 0.04$---statistically indistinguishable from zero---confirming that the enhancement is strictly mediated by vibronic resonance selection, not by trivial spectral narrowing of the pump pulse. 

This universality is particularly significant because it suggests that spectral bath engineering can be applied across a broad class of molecular aggregates where system-bath correlations dominate the transport dynamics. Whether in natural complexes or synthetic molecular wires, the principle of bypassing dissipative bath modes through resonant state preparation remains valid. The survival of this enhancement in a minimal three-site model (\cref{SI-sec:3site}) confirms that the effect is a fundamental property of \textbf{vibronically-coupled aggregates}: any molecular system where the electronic energy gap matches an underdamped vibrational mode ($\Delta E \approx \hbar\omega_{\mathrm{vib}}$) is a candidate for spectral coherence protection. This establishes a "top-down" control strategy that complements "bottom-up" chemical synthesis of light-harvesting materials.

\textbf{Environmental Robustness.}
We define the relative enhancement $\eta$ which quantifies the fractional gain from spectral filtering (\cref{eq:eta_def}). The enhancement persists across physiologically relevant conditions (\cref{fig:robustness}). Maximum enhancement occurs at \SIrange{285}{300}{\kelvin}, coinciding with the operating temperature of green sulfur bacteria. At \SI{295}{\kelvin}, $\eta = 0.25$; at \SI{310}{\kelvin}, $\eta$ decreases to $0.18$ but remains statistically significant ($p < 0.01$).

\begin{figure}[t]
\centering
\includegraphics[width=\columnwidth]{ETR_Under_Environmental_Effects.pdf}
\caption{\textbf{Environmental robustness.} \textbf{(a)}~Temperature dependence of $\eta(T)$, showing optimal enhancement at \SIrange{285}{300}{\kelvin}. \textbf{(b)}~Histogram of $\eta$ over \num{100} static disorder realizations ($\sigma = \SI{50}{\per\centi\meter}$): $\expval{\eta} = 0.20 \pm 0.04$. Error bars: \SI{95}{\percent} confidence intervals.}
\label{fig:robustness}
\end{figure}

Under static disorder ($\sigma = \SI{50}{\per\centi\meter}$), the $N=100$ production ensemble-averaged enhancement is $\expval{\eta} = 0.20(4)$ with a coefficient of variation of \SI{20}{\percent}. A comprehensive convergence and physical validation suite confirmed computational accuracy: \SI{1.8}{\percent} relative deviation from HEOM steady-state populations, trace conservation $< \num{5e-13}$, and density matrix eigenvalues bounded above $\num{-1e-14}$ (see SI Section~S3). The PT-HOPS method is formally exact; utilizing the globally synchronized production parameters ($L=8, K=2$), the automated convergence pipeline enforces strict hierarchy solver execution---failing automatically if lower-level approximations are invoked---and guarantees both trace preservation and positivity throughout the integration.

\textbf{Experimental Feasibility and Limitations.}
The proposed spectral filtering can be realized using programmable supercontinuum light sources with acousto-optic tuners,\cite{Gunthardt2023} or via commercial \SIlist{750;820}{\nano\meter} thin-film interference filters. We propose a specific 2DES implementation:\cite{Sebastian2024} Pump pulses (both $\vb{k}_1$ and $\vb{k}_2$) are spectrally shaped via a spatial light modulator (SLM) in a $4f$-shaper to match the dual-band transmission profile, while the probe pulse remains broadband.
The predicted coherence extension from \SI{280}{\femto\second} to \SI{420}{\femto\second} is well within the measurement resolution of current $\leq$\SI{50}{\femto\second} setups.\cite{Engel2007} Experimentally, this would manifest as prolonged cross-peak oscillations in the $\ket{1}$--$\ket{3}$ region of the 2D spectrum, accompanied by narrower anti-diagonal linewidths—a direct signature of reduced pure dephasing. 

Several approximations constrain our current analysis. The model is restricted to the single-exciton manifold and assumes site-independent baths with a phenomenological trapping rate. While we have ruled out alternative explanations such as increased photon flux (filtering reduces integrated flux) or vibrational coherence artifacts (electronic coherence dominates the $\gtrsim$\SI{300}{\femto\second} timescales), a full treatment of the pulse-shaping process including non-impulsive effects may reveal additional fine-tuning requirements (see SI Section~S11). We also note that in its natural state, FMO serves as a linker between the chlorosome baseplate and the reaction center and is not directly exposed to broadband sunlight.\cite{Blankenship2011} The proposed spectral control is therefore most directly relevant to laser-based 2DES experiments and to artificial bio-inspired light-harvesting architectures (such as our agrivoltaic concept), rather than asserting that this precise coherent excitation occurs natively in the biological setting.

\textbf{Conclusion and Outlook.}
PT-HOPS/SBD simulations show that quantum control via selective vibronic excitation---shaping the impulsive pump spectrum to align with vibronic resonances---extends excitonic coherence lifetimes by \SIrange{20}{50}{\percent} and improves forward transfer yields by \SI{25}{\percent} at room temperature. This mechanism relies on suppressing dephasing within the polaron frame by bypassing off-resonant thermal modes. The stability of this effect against disorder and temperature variations suggests potential for experimental realization.

Beyond natural photosynthesis, selective vibronic excitation offers a design path for artificial quantum materials. Synthetic systems, such as molecular J-aggregates, metal-organic frameworks (MOFs), and organic photovoltaics (OPVs), could be co-designed with photonic environments to shape the local vacuum and driving fields. In supramolecular J-aggregates, where the transition dipole moment is sensitive to the chromophore arrangement, targeted spectral excitation could suppress phonon-induced scattering channels and maintain long-range coherence in disordered environments.

Similarly, in MOFs, chemically tuning the vibrational modes of organic linkers provides a way to incorporate vibronic resonances into the material architecture. Combined with spectral filter technology, this approach may enable coherence-protected quantum channels for energy and information transport. The 2DES validation protocol proposed here provides a path for testing these principles. 

Ultimately, the spectral structure of the driving field acts as a controllable parameter for navigating the energy landscapes of condensed-phase quantum systems. In the context of sustainable engineering, this enables the definition of a \textbf{Symbiotic Power Conversion Efficiency (SPCE)}, where the dual-band filter is optimized not just for excitonic yield in the photovoltaic layer but for the simultaneous maintenance of the relative Electron Transport Rate (ETR) in underlying crops. This establishes quantum-coherent spectral engineering as a foundational technology for sustainable, symbiotic agrivoltaic ecosystems that address the energy--water--food nexus through fundamental physics.

%==============================================================================
% ACKNOWLEDGMENTS
%==============================================================================

\begin{acknowledgement}
This work was supported by the University of Yaoundé~I and the University of Douala. Computational resources were provided by departmental computing facilities. The MesoHOPS library was used for all PT-HOPS simulations.
\end{acknowledgement}

%==============================================================================
% DATA AVAILABILITY
%==============================================================================

\textbf{Data availability.} Simulation parameters and analysis scripts are available from the corresponding author upon reasonable request. The PT-HOPS/SBD implementation is based on the open-source MesoHOPS library.

%==============================================================================
% SUPPORTING INFORMATION
%==============================================================================

\begin{suppinfo}
PT-HOPS/SBD formalism derivation, FMO Hamiltonian parameters (site energies, electronic couplings, spectral density parameters), 12-test validation suite (hierarchy depth convergence, trace conservation, eigenvalue positivity, HEOM benchmark, high-temperature limit), convergence analysis, and computational cost breakdown.
\end{suppinfo}

%==============================================================================
% REFERENCES
%==============================================================================

\bibliography{references}

\end{document}
